% !TeX encoding = UTF-8
\section*{本章参考文献与阅读衔接}
\phantomsection
\addcontentsline{toc}{section}{本章参考文献与阅读衔接}

本章主线依据 \cite{yosida-zh,yosida-en}。原书第零章末的参考条目已核对：中译本与英文版均位于印刷第 22 页，即两份本地 PDF 的第 37 页。它按主题列出拓扑空间的三组参考作者和测度空间的两组参考作者，没有另列集合论与线性空间书目。原书的作者后缀“[1]”指各作者在全书参考书目中的第一条，\emph{不是本稿的参考文献编号}。对应关系如下。

\begin{enumerate}
 \item \textbf{P. Alexandroff--H. Hopf [1]}：\emph{Topologie}, Vol.~I, Springer, 1935，见 \cite{alexandroff-hopf}。英文原书总书目印刷第 469 页（PDF 第 484 页）使用拼写 Alexandrov；章末及中译本使用 Alexandroff，均指同一作者。
 \item \textbf{N. Bourbaki [1]}：\emph{Topologie G\'en\'erale}，Hermann，原书记录年份 1940--1942，见 \cite{bourbaki-topology}。核对位置为英文原书总书目印刷第 470 页（PDF 第 485 页）。这是分册出版的书目记录；不可把现代合订本的重印年份当作原书所引年份。
 \item \textbf{J. L. Kelley [1]}：\emph{General Topology}, D.~Van Nostrand, 1955，见 \cite{kelley-topology}。核对位置为英文原书总书目印刷第 475 页（PDF 第 490 页）。Springer 的 GTM 27 版为后来的再版，书目中另作说明。
 \item \textbf{P. R. Halmos [1]}：\emph{Measure Theory}, D.~Van Nostrand, 1950，见 \cite{halmos-measure}。核对位置为英文原书总书目印刷第 473 页（PDF 第 488 页）。
 \item \textbf{S. Saks [1]}：\emph{Theory of the Integral}, Warszawa, 1937，见 \cite{saks-integral}。核对位置为英文原书总书目印刷第 481 页（PDF 第 496 页）。
\end{enumerate}

以上五条均已并入书后的统一参考文献表，正文不再维护另一套编号。这里记录的是吉田原书的参考指引；本次实际核对了原书章末与总书目的对应条目，未把这五部参考书的全文标为已经逐页查阅。

阅读衔接上，第 \ref{sec:yosida-00-01} 节的 Zorn 引理用于代数基及延拓构造；第 \ref{sec:yosida-00-02} 节的紧性、类别定理与连续函数逼近分别服务于弱拓扑、Baire 方法和表示定理；第 \ref{sec:yosida-00-03} 节的积分工具为 $L^p$ 空间、卷积及积分算子准备基础；第 \ref{sec:yosida-00-04} 节的核、值域和商空间则在后续算子论中反复出现。第一章从这些基础出发，以半范数建立拓扑向量空间。
